% Background and Motivation: concepts, limitations, evidence, and challenges.
\section{Background and Motivation}
\label{sec:background-motivation}

\subsection{Background}
\label{sec:motivation-background}

\paragraph{Multi-agent execution.}
Many LLM applications rely on multi-agent collaboration to complete
complex tasks. These applications divide a user task into subtasks and
assign them to several agents for parallel execution. Each agent maintains
its own conversation, advancing its task through LLM requests and tool
invocations: the model generates responses or initiates tool calls, and
the returned results become input to subsequent
requests~\cite{wu2023autogen}. A \emph{request} is a single model call;
an \emph{agent} can initiate multiple requests to complete a subtask;
a \emph{workflow} connects these tasks through serial dependencies,
parallel branches and joins. The application orchestrates the workflow,
while the inference service executes its model requests.

Figure~\ref{fig:agent-workflow-examples} illustrates this structure.
In code generation, the main agent delegates work on a CSV-export feature
to API, UI, test-authoring and documentation agents. The branches execute
in parallel, alternating model requests with file reads, edits or checks.
The main agent then integrates the changes and runs tests. In deep
research, the lead agent delegates technical, cost, security and supply
investigations, then combines the collected evidence into a
report~\cite{anthropic2025multiagent}. The number of requests and tool
calls required by each branch may vary. Therefore, after a model request
completes, its agent and the entire workflow may still be unfinished.

\paragraph{Workflow completion time.}
Let $w$ identify a workflow, $a_w$ denote its arrival time at the system,
and $c_w$ denote the time when all its required outputs are ready.
We define \emph{workflow completion time} (WCT) as the elapsed time from
arrival to completion, $\mathrm{WCT}_w=c_w-a_w$.
It includes model execution, tool execution,
queueing and dependency waits. In the coding example, integration waits
for all required changes. A delay on the last unfinished branch postpones
integration, while an earlier-finishing branch may absorb the same delay
before the join. Request latency affects WCT through these dependencies.

\begin{figure}[t]
  \centering
  \includegraphics[width=\columnwidth]{figures/agent-workflow-examples.pdf}
  \caption{Illustrative multi-agent workflows.
    Each row is one agent: $R_i$ denotes its $i$th LLM request, with
    intervening tool results carried into the next request.
    Branches have different request counts and no shared round barriers.
    The coordinator resumes after their results join.
    Positions and lengths do not encode execution time.}
  \label{fig:agent-workflow-examples}
  \Description{Two stacked left-to-right workflows. A main agent assigns
    CSV-export API, UI, test-authoring and documentation tasks in parallel,
    then integrates their outputs and runs tests. Their request counts
    are four, three, three and two, respectively. A lead researcher assigns
    technical, cost, safety and supply investigations, with four, two,
    three and two requests, before writing a report. Each branch explicitly
    unfolds its requests and intervening tools. Starts and ends are
    staggered; all edges move forward, without loops or round barriers.}
\end{figure}

\paragraph{KV reuse in multi-agent workflows.}
The model requests on these branches repeatedly consume context from
earlier steps. Coding agents carry source files, edits and test outputs
into later requests; research agents accumulate retrieved documents,
intermediate findings and tool results. Processing this growing history
from scratch on every request would repeat work already performed.
During \emph{prefill}, an LLM computes per-layer keys and values (KV)
for input tokens; \emph{decoding} extends this state as output tokens
are generated. Retaining this KV allows a subsequent request with a
matching token prefix to compute only the new suffix, provided the model
and position semantics are compatible~\cite{gao2024cachedattention,zheng2024sglang}.

This reuse occurs both within and across agents. Successive requests
from one agent can reuse KV for their shared conversation history.
Different agents can also inherit a common prompt before adding their
own tasks. For example, the API, UI and test agents in
Figure~\ref{fig:agent-workflow-examples} can share KV for common project
instructions and repository context when these form the same initial
token sequence. Different role prompts or context edits can shorten
the reusable prefix~\cite{zheng2024sglang}.
The KV cache therefore preserves computation across requests, including
while an agent executes tools or waits for another branch.

\paragraph{Failures and KV reconstruction.}
These reuse benefits depend on the cached state remaining available.
In large-scale LLM deployments, hardware faults and software crashes can
interrupt serving workers and make their GPU-resident KV
unavailable~\cite{cao2026lumen,jayakody2026ghostserve}.
For production-cluster context, Figure~\ref{fig:gpu-server-availability}
shows the public InfiniteHBD training-cluster trace: 584 recorded outages
affect 231 of 400 GPU servers over approximately 349 days, with reconstructed
availability falling to 365 servers~\cite{shou2025infinitehbd}.
The trace includes faults and operational interruptions, as detailed in
the caption; it characterizes infrastructure availability rather than
serving recovery latency.

When no reusable KV survives, a replacement worker must reconstruct it
by replaying the retained tokens, even if it already has the model weights
and the application has preserved the conversation
history~\cite{strati2024dejavu}. For a long-running agent, this means
reprocessing the context accumulated across earlier requests rather than
only the latest input~\cite{jayakody2026ghostserve}.
One failed worker can also interrupt several co-located agents, requiring
surviving workers to reconstruct multiple histories while handling normal
traffic~\cite{cao2026lumen}. The resulting delays propagate through the
workflow's dependencies: reconstruction on the last unfinished branch
postpones its join and can increase WCT.

\begin{figure}[t]
  \centering
  \includegraphics[width=\columnwidth]{figures/gpu-server-availability.pdf}
  \caption{GPU-server availability reconstructed from InfiniteHBD's
    production-training trace~\cite{shou2025infinitehbd}: 400 servers minus
    those within recorded outage intervals. Records include GPU, network,
    power/cooling, and software faults; unknown causes, tests,
    troubleshooting, and changes are also retained. Overlapping intervals
    per server are merged. Availability is measured at the node level.}
  \label{fig:gpu-server-availability}
  \Description{A step curve over approximately 349 days shows 365 to 400
    GPU servers outside active recorded fault intervals in a fixed
    400-server cohort. The dashed line marks all 400 servers.}
\end{figure}

\paragraph{KV backup across storage tiers.}
To reduce reconstruction work, a \emph{KV checkpoint} retains an already
computed prefix. If it survives the failure, recovery retrieves this
prefix and replays only the unsaved
tail~\cite{strati2024dejavu,jayakody2026ghostserve}.
KV can be retained in GPU HBM, host DRAM, or shared
storage~\cite{gao2024cachedattention,qin2025mooncake}.
These tiers differ in both storage cost and recovery latency.
HBM offers fast access but has high capacity cost and competes with
foreground inference for GPU memory. DRAM and storage-backed shared
services can provide cheaper capacity, with additional retrieval latency.

Figure~\ref{fig:motivation-recompute} illustrates the latency difference
between reconstructing a history and reusing retained KV.
At 1,010,000 tokens, full prefill takes 571.21\,s even with sparse
attention; HBM, DRAM and shared-storage reuse take 1.305, 2.663 and
10.156\,s, respectively. The history occupies 57.92\,GB across two ranks.
All paths end when the same input's KV is ready: reuse restores the
preceding history, if needed, and computes one new token.
HBM is a resident-state reference. These measurements quantify KV
readiness, excluding checkpoint creation, failure detection and routing.

\begin{figure}[t]
  \centering
  \includegraphics[width=\columnwidth]{figures/recovery-cost.pdf}
  \caption{Measured KV reconstruction and reuse costs for
    Qwen2.5-7B-Instruct-1M on two A100 PCIe 40\,GB GPUs
    (TP2, BF16, batch one), with sparse attention above 32K.
    Full prefill starts from empty KV; reuse starts from $(C-1)$ cached
    tokens and computes one. KV reconstruction time includes any
    retrieval and computation until all $C$ tokens' KV is ready.
    DRAM/shared-storage I/O is serialized with computation.
    HBM/DRAM: median and min--max of three timings;
    shared storage/full prefill: one timing per length.
    Protocol: Appendix~\ref{sec:motivation-protocol}.
    K~=~1,024; M~=~1,000,000.}
  \label{fig:motivation-recompute}
  \Description{Four measured KV reconstruction-time curves for Qwen2.5-7B-Instruct-1M
    from 8K to 1,010,000 tokens: blue circles for resident HBM reuse,
    orange squares for DRAM reuse, purple diamonds for shared-storage
    reuse, and red triangles for full prefill. Both axes are logarithmic.
    At the longest context, the four paths take 1.305, 2.663, 10.156 and
    571.21 seconds, respectively.}
\end{figure}

\subsection{Motivation and Challenges}

The storage tiers above provide different ways to retain KV and avoid
reconstruction. Their capacity costs and access latencies leave a further
question: which state should be protected, and where, to reduce
failure-induced increases in WCT under a limited storage budget?
Fault-tolerant serving and workflow-aware KV management approach this
question from different directions. The following limitations explain
why their existing priorities do not directly resolve it.

\begin{figure*}[t]
  \centering
  \captionsetup[subfigure]{font=normalsize,labelfont=normalfont,
    textfont=normalfont,justification=centering,singlelinecheck=false}
  \begin{subfigure}[t]{2.48in}
    \vspace{0pt}
    \begin{minipage}[c][1.65in][c]{\linewidth}
      \centering
      \includegraphics[width=\linewidth]{figures/workflow.pdf}
    \end{minipage}
    \caption{Four-branch workflow}
    \label{fig:motivation-workflow}
  \end{subfigure}\hfill
  \begin{subfigure}[t]{1.86in}
    \vspace{0pt}
    \begin{minipage}[c][1.65in][c]{\linewidth}
      \centering
      \includegraphics[width=\linewidth]{figures/workflow-gaps.pdf}
    \end{minipage}
    \caption{Workflow-maximum gaps}
    \label{fig:motivation-branch-gap}
  \end{subfigure}\hfill
  \begin{subfigure}[t]{2.52in}
    \vspace{0pt}
    \begin{minipage}[c][1.65in][c]{\linewidth}
      \centering
      \includegraphics[width=\linewidth]{figures/recovery-exposure.pdf}
    \end{minipage}
    \caption{Recovery cost and WCT}
    \label{fig:motivation-recovery-exposure}
  \end{subfigure}
  \caption{Workflow timing and conditional recovery impact.
    (a) Four trace-derived branches with schematic Main stages A0/A5.
    Bars preserve timing; dashed tails assume an all-branch join.
    (b) ECDF of maximum branch-duration gaps per workflow: 121 workflows
    with temporally inferred fan-in; these are not failure-time slack.
    (c) Measured Prefill costs at the four branches' trace input-size
    proxies, using Figure~\ref{fig:motivation-recompute}'s configuration.
    Bars/whiskers show median/min--max over three timings after warmup.
    Hatching denotes absorbed delay; solid fill is the calculated WCT
    increase under an all-branch join with other execution fixed.
    Each case is a single-branch counterfactual, not a fault-injected WCT
    measurement. Protocol: Appendix~\ref{sec:motivation-protocol}.}

  \label{fig:motivation}
  \Description{Left: a four-branch workflow with A2 completing last at
    151.6 seconds. Middle: an ECDF of workflow-maximum branch-duration
    differences with a median of 138 seconds and a 90th percentile of
    1173 seconds. Right: A1 through A4 have measured median Prefill costs
    of 33.2, 32.6, 35.4 and 39.9 seconds. Applying each cost separately
    as an added branch pause gives calculated WCT increases of 0.0, 32.6,
    16.3 and 18.0 seconds. Hatching identifies the absorbed portion;
    whiskers show timing minima and maxima, not confidence intervals.}
\end{figure*}

\paragraph{Limitation 1: Fault tolerance without workflow-aware priorities.}
Fault-tolerant serving systems preserve state and coordinate recovery
for interrupted requests. D\'ej\`aVu streams KV checkpoints, and GhostServe
uses host-memory parity to reduce reconstruction
work~\cite{strati2024dejavu,jayakody2026ghostserve}.
LUMEN further coordinates checkpoint placement, recovery routing and
capacity restoration according to worker load~\cite{cao2026lumen}.
These mechanisms address state loss and recovery contention, but do not
use multi-agent dependencies to determine which requests' faster
recovery would advance workflow completion. Request recovery cost and
worker load alone cannot express this benefit.

\label{sec:motivation-workflow}
Figure~\ref{fig:motivation}(\subref{fig:motivation-workflow}) illustrates
why workflow position matters. A2 finishes last, 151.6\,s after the fork,
whereas A1 finishes at 60.4\,s. If the following Main stage requires all
four outputs, A1 can absorb 91.2\,s of additional delay before postponing
the join, whereas any delay on A2 postpones it. We call the additional
delay a branch can absorb without postponing the required output its
\emph{slack}.

Such timing differences occur across the retained CC/Weka
traces~\cite{semianalysis2026cctraces}. From 393 filtered Claude Code
sessions, a temporal fan-in proxy retains 286 overlapping subagent
groups in 121 workflows: all branches complete, and a Main request
follows the last completion; Main activity during the branches is allowed.
Figure~\ref{fig:motivation}(\subref{fig:motivation-branch-gap}) reports
each workflow's maximum longest-minus-shortest branch-duration gap,
with a median of 138\,s and a 90th percentile of 1,173\,s.
The trace lacks explicit result-consumption edges, so these differences
neither verify all-branch joins nor measure failure-time slack.
The four-branch example is one retained group, not a representative
workflow-maximum gap.

To connect reconstruction cost to workflow delay, we measure full
prefill at the four active requests' trace input-size proxies, using
Figure~\ref{fig:motivation-recompute}'s Qwen2.5-7B-Instruct-1M
configuration with sparse attention~\cite{qwen2025qwen25oneM,jiang2024minference}.
Figure~\ref{fig:motivation}(\subref{fig:motivation-recovery-exposure})
then inserts each measured cost $R_i$ separately as a pause on branch
$i$ at 29.5\,s after the fork. With offline slack $\sigma_i$, an assumed
all-branch join, and all other execution and queues fixed, the WCT
increment is
\[
  \Delta\mathrm{WCT}_i=\max(0,R_i-\sigma_i).
\]
A1's 33.2\,s reconstruction fits within its 91.2\,s slack and leaves WCT
unchanged. A2's slightly smaller cost, 32.6\,s, is fully exposed because
its slack is zero. Thus prioritizing the more expensive reconstruction
can favor a branch with less immediate effect on completion.
Only the Prefill costs are measured; WCT increments are conditional
calculations using synthetic inputs and trace timing, not fault-injected
workflow measurements. Appendix~\ref{sec:motivation-protocol} gives the
input construction, timing samples and mapping protocol.

\paragraph{Challenge 1: Estimating workflow recovery urgency online.}
Protection decisions must reflect how a recovery delay propagates to
required workflow outputs. The example obtains slack from completed
traces, whereas an online system must act before the remaining model
and tool durations are known. The branch that gates completion can
change as execution progresses. A worker failure can also interrupt
several agents whose recoveries share resources, so isolated WCT
increments cannot simply be added. The challenge is to derive and
update useful protection priorities from partial workflow progress
under uncertain remaining execution and recovery load.

\paragraph{Limitation 2: Workflow-aware caching without a recovery objective.}
Workflow-aware KV management uses application structure to improve
normal execution. KVFlow uses upcoming agent invocations to guide prefix
retention and prefetching~\cite{pan2025kvflow}; TOPAS jointly schedules
prefix residency and request admission using remaining workflow paths,
reuse benefits and movement or preemption costs~\cite{ni2026topas}.
These systems already exploit workflow information, but their cache
decisions do not explicitly optimize protection against serving-worker
failures: which state must survive, where to retain it, and how quickly
it should become available to a replacement worker.

The distinction changes the value of retained state. For example,
KVFlow gives varying suffixes the highest eviction priority while
protecting reusable fixed prompts~\cite{pan2025kvflow}. A suffix with
little expected reuse can nevertheless be needed to resume an
interrupted execution. Likewise, keeping a prefix in the executing
worker's HBM accelerates normal access but does not by itself preserve
that state when the worker fails. A cache copy can support recovery if
it survives; its usefulness for future calls alone does not determine
its value as protection against a particular failure.

\label{sec:motivation-long-context}
The tier costs in Figure~\ref{fig:motivation-recompute} make the resulting
allocation tradeoff concrete. Retaining KV in a slower tier already
avoids much of the full-prefill cost; faster tiers reduce the remaining
retrieval delay at a higher capacity cost. Whether that reduction
advances workflow completion depends on the recovering branch's slack.
Conversely, a branch with ample slack may still reuse its history during
normal execution, so low recovery urgency does not justify deleting its
last reusable copy. These are distinct reasons to retain and accelerate
state, even when both use the same storage hierarchy.

\paragraph{Challenge 2: Allocating KV protection across storage tiers.}
A protection policy must decide both which histories remain recoverable
from saved KV and which merit faster access, accounting for workflow
recovery urgency, normal reuse and tier capacity costs. These demands
change as agents accumulate context, invoke tools and release dependent
work. The policy must preserve useful state outside the protected
failure boundary while limiting the capacity and copying work taken
from foreground service. Simply adding checkpoints to a reuse-oriented
cache does not determine these priorities; managing the two objectives
requires a coordinated retention and placement policy.

\label{sec:motivation-goals}
AegisKV addresses this gap by using multi-agent workflow progress to
guide KV protection and by allocating limited tiered storage to recovery
and normal reuse. The next section develops the observations and
mechanisms that support these decisions.
